\section{Roteiro Prático 03 — Atividade Prática Autônoma: Calculadora de Pedidos \& Carrinho de Compras}

% =====================================================
% OBJETIVOS DA ATIVIDADE
% =====================================================

\begin{boxObjetivos}
\textbf{Objetivos da Atividade (Atividade Guiada de 2 Semanas — Calculadora de Pedidos)}

Ao final desta atividade prática autônoma, o estudante será capaz de:

\begin{itemize}
\item Consolidação prática e aplicação integrada do \textbf{Bootstrap 5} e do \textbf{jQuery} em uma aplicação web funcional;
\item Ler valores de campos de formulário (\texttt{<select>}, \texttt{<input type="number">}, \texttt{<input type="checkbox">}) e convertê-los com segurança para tipos numéricos utilizando \texttt{parseFloat()} e \texttt{parseInt()};
\item Tratar valores vazios ou não numéricos (\texttt{NaN}) utilizando o operador de fallback \texttt{|| 0};
\item Compreender a diferença e aplicar os eventos corretos do navegador: \texttt{.on('input')} para digitação em tempo real, \texttt{.on('change')} para alteração de seleção/switches e \texttt{.on('click')} para botões;
\item Realizar o cálculo automático do subtotal, taxa de entrega, descontos por cupom e total geral em tempo real, sem necessidade de recarregar a página;
\item Formatar valores monetários para o padrão brasileiro (\texttt{R\$ X,XX}) utilizando a função \texttt{.toFixed(2).replace('.', ',')};
\item Persistir e restaurar o rascunho do pedido no navegador utilizando o \textbf{LocalStorage} (\texttt{setItem}, \texttt{getItem}, \texttt{JSON.stringify}, \texttt{JSON.parse});
\item Organizar o projeto em arquivo externo \texttt{app.js} e publicar a aplicação no \textbf{GitHub Pages}.
\end{itemize}
\end{boxObjetivos}

% =====================================================
% INTRODUÇÃO \& CONTEXTUALIZAÇÃO
% =====================================================

\subsection{Introdução e Contextualização da Atividade}

Prezados estudantes da disciplina de \textbf{Programação para a Internet 2 (ProgWeb 2)}!

Durante as próximas 2 semanas, enquanto o professor estará em afastamento de saúde, vocês realizarão uma atividade prática autônoma guiada para consolidar os conceitos aprendidos sobre \textbf{Bootstrap 5}, \textbf{jQuery}, \textbf{Cálculos Numéricos no Cliente} e \textbf{Persistência no LocalStorage}.

O tema proposto é o desenvolvimento da \textbf{Lanchonete Express — Calculadora de Pedidos \& Carrinho de Compras}. A aplicação permitirá que o cliente escolha seu lanche principal, selecione adicionais opcionais (queijo, bacon, batata frita, bebida), ajuste a quantidade, escolha o tipo de entrega, aplique um cupom de desconto e veja o \textbf{valor total a pagar atualizado instantaneamente na tela}.

% =====================================================
% FUNDAMENTAÇÃO TEÓRICA E DICAS DE CÓDIGO
% =====================================================

\subsection{Fundamentação Teórica: Como Ler e Calcular Valores de Campos com jQuery}

Para realizar cálculos financeiros em uma página web, é necessário seguir três etapas indispensáveis:

\subsubsection{1. Leitura e Conversão de Texto para Número (\texttt{parseFloat} e \texttt{parseInt})}

No HTML, todos os valores obtidos através de \texttt{\textdollar(seletor).val()} são retornados como \textbf{texto (string)}. Se tentarmos somar textos em JavaScript (ex: \texttt{"18.00" + "4.00"}), o resultado será uma concatenação (\texttt{"18.004.00"}), e não uma soma numérica.

Por isso, devemos converter o valor lido:
\begin{itemize}
\item \texttt{parseFloat(string)}: Converte o texto para número decimal (usado para preços);
\item \texttt{parseInt(string)}: Converte o texto para número inteiro (usado para quantidades);
\item \textbf{Tratamento de NaN (Not a Number):} Caso o campo esteja vazio, o \texttt{parseFloat()} retornará \texttt{NaN}. Para evitar que os cálculos quebrem, usamos a instrução \texttt{|| 0}.
\end{itemize}

\begin{boxCodigo}
\textbf{Exemplo de leitura segura com fallback zero:}

\small
\begin{verbatim}
const precoLanche = parseFloat($('#select-lanche').val()) || 0;
const quantidade = parseInt($('#input-qtd').val()) || 1;
\end{verbatim}
\end{boxCodigo}

\subsubsection{2. Soma de Múltiplos Checkboxes Marcados com \texttt{.each()}}

Quando a aplicação possui múltiplos adicionais (ex: switches de queijo, bacon, batata frita), podemos percorrer apenas os checkboxes que o usuário marcou usando a combinação do seletor \texttt{:checked} com o método \texttt{.each()}:

\begin{boxCodigo}
\textbf{Varredura de adicionais selecionados:}

\small
\begin{verbatim}
let somaAdicionais = 0;

$('.check-adicional:checked').each(function() {
  somaAdicionais += parseFloat($(this).val()) || 0;
});
\end{verbatim}
\end{boxCodigo}

\subsubsection{3. Guia de Eventos do Navegador (\texttt{input}, \texttt{change} e \texttt{click})}

Para proporcionar uma excelente experiência de usuário (UX), a calculadora deve reagir instantaneamente a qualquer alteração realizada na tela. Escolha o evento adequado para cada tipo de elemento:

\begin{boxTabela}
\textbf{Guia de Eventos do Navegador em jQuery}

\begin{tabularx}{\linewidth}{|l|X|}
\hline
\textbf{Evento jQuery} & \textbf{Componente Indicado e Comportamento} \\ \hline
\texttt{.on('change')} & Indicado para \texttt{select}, \texttt{checkbox} e \texttt{switch}. Disparado quando o usuário altera a seleção de um dropdown ou marca/desmarca uma caixa. \\ \hline
\texttt{.on('input')} & Indicado para \texttt{input type="number"} ou \texttt{text}. Disparado instantaneamente a cada caractere digitado ou alterado no campo. \\ \hline
\texttt{.on('click')} & Indicado para botões. Disparado quando o usuário clica nos botões de incremento (+/-), aplicar cupom ou finalizar pedido. \\ \hline
\end{tabularx}
\end{boxTabela}

\subsubsection{4. Formatação Financeira em Real (\texttt{R\$})}

Para exibir valores numéricos formatados na moeda brasileira (\texttt{R\$ 25,50}), crie uma função auxiliar utilizando \texttt{.toFixed(2)} para fixar duas casas decimais e \texttt{.replace('.', ',')} para trocar o ponto pela vírgula:

\begin{boxCodigo}
\textbf{Função de Formatação de Moeda Brasileira:}

\small
\begin{verbatim}
function formatarMoeda(valor) {
  return 'R$ ' + valor.toFixed(2).replace('.', ',');
}

// Atualização da interface gráfica:
$('#total-geral').text(formatarMoeda(totalCalculado));
\end{verbatim}
\end{boxCodigo}

% =====================================================
% DICAS DE PAINEIS E LAYOUT BOOTSTRAP 5
% =====================================================

\subsection{Dicas de Estruturação Visual \& Componentes do Bootstrap 5}

Para que a aplicação tenha um visual profissional e limpo, recomendamos a utilização dos seguintes painéis e utilitários do \textbf{Bootstrap 5}:

\begin{enumerate}
\item \textbf{Grid de 2 Colunas (\texttt{.row} e \texttt{.col-lg-7} / \texttt{.col-lg-5}):} Divida a tela em dois painéis principais — a coluna da esquerda para o formulário de montagem do pedido e a coluna da direita para o resumo financeiro;
\item \textbf{Painel de Montagem (\texttt{.card} e \texttt{.card-body}):} Agrupe as opções do pedido em um cartão elevado com bordas suaves;
\item \textbf{Adicionais com Switches (\texttt{.form-check .form-switch}):} Utilize interruptores modernos para os adicionais opcionais (queijo, bacon, fritas), incluindo a etiqueta \texttt{.badge} com o preço;
\item \textbf{Input Groups para Quantidade (\texttt{.input-group}):} Combine os botões de decremento (\texttt{-}) e incremento (\texttt{+}) com o campo numérico no meio;
\item \textbf{Resumo Financeiro em Lista (\texttt{.list-group .list-group-flush}):} Exiba a discriminação do subtotal, taxa de entrega e desconto em uma lista estilizada;
\item \textbf{Destaque do Total Geral (\texttt{.alert} ou caixa personalizada):} Destaque o valor total a pagar com tamanho de fonte ampliado e fundo contrastante.
\end{enumerate}

% =====================================================
% PASSO A PASSO GUIADO DA ATIVIDADE
% =====================================================

\subsection{Roteiro Guiado Passo a Passo (Roteiro 3)}

Siga os passos abaixo para construir a aplicação na sua pasta de exercícios (\texttt{roteiro-3/exercicio-roteiro3/}):

\begin{boxAtividade}[{Passo 1: Estrutura HTML5 e Inclusão das CDNs}]
Crie o arquivo \texttt{index.html} e inclua as bibliotecas do \textbf{Bootstrap 5}, \textbf{Bootstrap Icons} e \textbf{jQuery 3.7.1} via CDN antes de fechar a tag \texttt{</body>}.

\begin{boxCodigo}
\small
\begin{verbatim}
<!-- Bootstrap 5 CSS -->
<link href="https://cdn.jsdelivr.net/npm/bootstrap@5.3.3/dist/css/bootstrap.min.css" 
      rel="stylesheet">
<!-- jQuery CDN -->
<script src="https://code.jquery.com/jquery-3.7.1.min.js"></script>
<!-- Arquivo da sua lógica JS -->
<script src="app.js"></script>
\end{verbatim}
\end{boxCodigo}
\end{boxAtividade}

\begin{boxAtividade}[{Passo 2: Formulário de Montagem do Pedido}]
Crie o menu de seleção do lanche principal atribuindo os preços no atributo \texttt{value} de cada \texttt{<option>}.

\begin{boxCodigo}
\small
\begin{verbatim}
<select class="form-select" id="select-lanche">
  <option value="0">-- Selecione um lanche --</option>
  <option value="18.00">X-Burger Clássico — R$ 18,00</option>
  <option value="22.00">X-Salada Especial — R$ 22,00</option>
  <option value="26.00">X-Bacon Supremo — R$ 26,00</option>
</select>
\end{verbatim}
\end{boxCodigo}
\end{boxAtividade}

\begin{boxAtividade}[{Passo 3: Lógica de Cálculo no app.js}]
No arquivo \texttt{app.js}, estruture a função \texttt{\textdollar(document).ready()} e crie a função principal \texttt{calcularTotal()}:

\begin{boxCodigo}
\small
\begin{verbatim}
$(document).ready(function() {

  function calcularTotal() {
    const precoLanche = parseFloat($('#select-lanche').val()) || 0;
    
    let somaAdicionais = 0;
    $('.check-adicional:checked').each(function() {
      somaAdicionais += parseFloat($(this).val()) || 0;
    });

    let qtd = parseInt($('#input-qtd').val()) || 1;
    const subtotal = (precoLanche + somaAdicionais) * qtd;
    const taxaEntrega = parseFloat($('#select-entrega').val()) || 0;

    const totalGeral = subtotal + taxaEntrega;
    
    $('#total-geral').text('R$ ' + totalGeral.toFixed(2).replace('.', ','));
  }

  // Escutadores em tempo real
  $('#select-lanche, #select-entrega, .check-adicional').on('change', calcularTotal);
  $('#input-qtd').on('input change', calcularTotal);

});
\end{verbatim}
\end{boxCodigo}
\end{boxAtividade}

\begin{boxAtividade}[{Passo 4: Persistência no LocalStorage}]
Implemente o salvamento automático do rascunho do pedido no \texttt{localStorage} ao clicar em "Finalizar Pedido" e a restauração automática ao pressionar \texttt{F5}:

\begin{boxCodigo}
\small
\begin{verbatim}
// Salvar no LocalStorage
$('#btn-finalizar').on('click', function() {
  const pedido = {
    lanche: $('#select-lanche').val(),
    qtd: $('#input-qtd').val(),
    entrega: $('#select-entrega').val()
  };
  localStorage.setItem('rascunho_pedido', JSON.stringify(pedido));
  alert('Pedido salvo no navegador!');
});
\end{verbatim}
\end{boxCodigo}
\end{boxAtividade}

% =====================================================
% CHECKLIST E ENTREGÁVEL DO ESTUDANTE
% =====================================================

\subsection{Checklist de Conclusão e Orientações de Entrega}

Antes de enviar sua atividade, verifique se você concluiu todos os itens da lista abaixo:

\begin{itemize}
\item[$\Box$] A página carrega corretamente sem erros no console do navegador (\texttt{F12});
\item[$\Box$] O layout é responsivo e utiliza os cartões e formulários do Bootstrap 5;
\item[$\Box$] O cálculo do valor total reage em tempo real a qualquer mudança nos selects, checkboxes ou quantidade;
\item[$\Box$] Os valores numéricos são devidamente formatados em Reais (\texttt{R\$ X,XX});
\item[$\Box$] O rascunho do pedido é salvo e restaurado via \texttt{localStorage};
\item[$\Box$] O código foi commitado no repositório Git pessoal e publicado via \textbf{GitHub Pages}.
\end{itemize}
